\subsection{محدودیت‌ها و تحلیل شکست}
تحلیل چندوجهی کارایی این ماژول نکات برجسته‌ای را آشکار می‌سازد:
\begin{itemize}
\item \textbf{پایداری همه‌جانبه با نمره قبولی کامل (PASS):} تمامی چهار مدل تحت آزمون موفق به کسب نرخ بالای دقت (بین \lr{84.63\%} تا \lr{89.11\%}) با همگرایی پایدار لاگرانژی شدند.
\item \textbf{تعادل قابل توجه دقت و مصرف حافظه:} مدل \lr{RoBERTa} با دقت \lr{89.11\%} و F1 معادل \lr{0.8966} و مصرف حافظه تنها \lr{333.44 MB}، گواهی بر قدرت همگرایی هم‌زمان هرس و کوانتیزاسیون است.
\item \textbf{چالش زمان کالیبراسیون و بار محاسباتی بهینه‌سازی:} علی‌رغم نتایج استنتاجی عالی، اجرای تکرارهای متناوب حل بسته و تصاویر غیرخطی نیازمند صرف زمان پردازشی اولیه و حافظه محاسباتی برای نگهداری متغیرهای دوگان ($\mathbf{U}_1, \mathbf{U}_2$) در طول فشرده‌سازی است. در صورتی که تعداد تکرارها کمتر از ۴ تنظیم شود، واگرایی دوگان رخ داده و خطای بازسازی به شدت افزایش می‌یابد.
\end{itemize}
